\documentclass[11pt,a4paper]{article}

% ──────────────────────────────────────────────────────────────────────────
% AXON-RFC-006 — Stateful EasyNet (Minimal Normative Core)
% Status: Draft v0.1 (pre-implementation main body)
% Build:  xelatex AXON-RFC-006-stateful-easynet.tex
% Companion: AXON-RFC-006-stateful-easynet.md (Appendix B pressure test)
% ──────────────────────────────────────────────────────────────────────────

\usepackage[margin=1in]{geometry}
\usepackage{amsmath,amssymb,amsthm}
\usepackage{xcolor}
\usepackage[colorlinks=true,linkcolor=blue!60!black,citecolor=blue!60!black,urlcolor=blue!60!black]{hyperref}
\usepackage{booktabs}
\usepackage{longtable}
\usepackage{array}
\usepackage{tabularx}
\usepackage{xltabular}
\newcolumntype{L}[1]{>{\raggedright\arraybackslash}p{#1}}
\usepackage{listings}
\usepackage{enumitem}
\usepackage{tcolorbox}
\usepackage{titlesec}
\usepackage{fancyhdr}
\usepackage{parskip}
\usepackage{xspace}

% ── Unicode + CJK support (xelatex) ───────────────────────────────────────
\usepackage{fontspec}
\usepackage{xeCJK}
\setCJKmainfont{PingFang SC}[
    BoldFont={PingFang SC Semibold},
    ItalicFont={PingFang SC Light}
]

% ── Page style ────────────────────────────────────────────────────────────
\pagestyle{fancy}
\fancyhf{}
\fancyhead[L]{\small AXON-RFC-006}
\fancyhead[R]{\small Stateful EasyNet --- Minimal Normative Core}
\fancyfoot[C]{\thepage}
\renewcommand{\headrulewidth}{0.4pt}

% ── Section styling ───────────────────────────────────────────────────────
\titleformat{\section}{\Large\bfseries\color{black!85}}{\thesection}{0.7em}{}
\titleformat{\subsection}{\large\bfseries\color{black!75}}{\thesubsection}{0.6em}{}
\titleformat{\subsubsection}{\normalsize\bfseries\color{black!70}}{\thesubsubsection}{0.5em}{}

% ── Tag colors ────────────────────────────────────────────────────────────
\definecolor{normativecolor}{RGB}{180,30,30}
\definecolor{derivedcolor}{RGB}{30,90,180}
\definecolor{deferredcolor}{RGB}{120,90,30}
\definecolor{codebg}{RGB}{245,245,248}
\definecolor{codeborder}{RGB}{210,210,220}

% ── Tag macros (parallel to AXIOM.tex's [axiom]/[derived]/[deferred]) ─────
\newcommand{\normsv}{\textcolor{normativecolor}{\textbf{[normative]}}\xspace}
\newcommand{\derived}{\textcolor{derivedcolor}{\textbf{[derived]}}\xspace}
\newcommand{\deferred}{\textcolor{deferredcolor}{\textbf{[deferred]}}\xspace}

% ── Code listing style ────────────────────────────────────────────────────
\lstdefinestyle{ealstyle}{
    backgroundcolor=\color{codebg},
    basicstyle=\ttfamily\small,
    frame=single,
    rulecolor=\color{codeborder},
    breaklines=true,
    showstringspaces=false,
    keywordstyle=\color{blue!70!black}\bfseries,
    commentstyle=\color{green!40!black}\itshape,
    stringstyle=\color{red!60!black},
    morekeywords={state_type,state_key,canonical,operational,transition,receipt,view,principal,agent,owner_signature,attempt_id,transition_id},
    columns=fullflexible,
    keepspaces=true,
    aboveskip=0.6em,
    belowskip=0.6em,
}
\lstset{style=ealstyle}

% ── Theorem-like environments ─────────────────────────────────────────────
\newtheorem{theorem}{Theorem}
\newtheorem{lemma}[theorem]{Lemma}
\newtheorem{corollary}[theorem]{Corollary}
\theoremstyle{definition}
\newtheorem{definition}{Definition}

% Invariant box accepts an optional tag in #1 (e.g. SO-INV-1) and
% renders it inside the title bar so we don't have to push it into
% pgfkeys (which would treat the tag as an unknown key).
\newtcolorbox{invariant}[1][Invariant]{
    colback=red!4!white,
    colframe=red!50!black,
    fonttitle=\bfseries,
    title={Invariant --- #1},
    sharp corners=south,
    boxrule=0.6pt,
}
\newtcolorbox{notebox}[1][]{
    colback=blue!4!white,
    colframe=blue!50!black,
    fonttitle=\bfseries,
    title={Note},
    sharp corners=south,
    boxrule=0.6pt,
    #1
}
\newtcolorbox{warningbox}[1][]{
    colback=orange!6!white,
    colframe=orange!70!black,
    fonttitle=\bfseries,
    title={Transitional misalignment},
    sharp corners=south,
    boxrule=0.6pt,
    #1
}
\newtcolorbox{decisionbox}[1][]{
    colback=purple!4!white,
    colframe=purple!60!black,
    fonttitle=\bfseries,
    title={Open decision},
    sharp corners=south,
    boxrule=0.6pt,
    #1
}

% ── Title ─────────────────────────────────────────────────────────────────
\title{
    \vspace{-2em}
    \textbf{\Large AXON-RFC-006: Stateful EasyNet}\\[0.3em]
    \large Minimal Normative Core --- the protocol primitives a stateful
    agent network requires
}
\author{
    Silan Hu \\ \small National University of Singapore
}
\date{Draft v0.1, \today}

% ──────────────────────────────────────────────────────────────────────────
\begin{document}
\maketitle

\begin{abstract}
\noindent
This document defines the minimal protocol surface a stateful EasyNet
deployment must implement. It deliberately limits itself to seven
sections covering the primitives that Phase~1 implementation work
requires --- state object identity, canonical state hashing, receipt
class separation, transition ability schema, authority and
delegation, view projection, and receipt indexing by state key ---
and explicitly defers everything else (full theorems, hub
deployment, web-app and pages reference machines) to companion
artefacts. The companion document \texttt{AXON-RFC-006-stateful-easynet.md}
contains a pressure-test appendix (Appendix~B,
\texttt{easynet.web\_app}) whose thirteen
invariants this main body is engineered to satisfy. The relationship
between the two documents is intentional: the appendix was written
\emph{first}, as the hardest concrete state object the protocol must
keep consistent; this main body then commits the minimal vocabulary
and rules under which the appendix's invariants hold. Section~8 lists
a coverage matrix proving that every appendix invariant has at least
one normative anchor here, and that every locally-introduced
invariant in this document (SO/CSH/RC/TA/VP-INV-*) is self-contained.
The work this document directly authorises is Phase~1: schema-level
\texttt{canonical:bool} marking, the two-class receipt split,
\texttt{principal/} caller URIs, a content-as-blob boundary, and a
state-key-keyed receipt index.
\end{abstract}

\tableofcontents
\vspace{1em}
\hrule
\vspace{1em}

% ──────────────────────────────────────────────────────────────────────────
\section{Reading guide}
\label{sec:reading-guide}

Every substantive statement in this document carries one of three
tags, identical in role to the tags used in the Axon Invocation
Axiom paper:

\begin{itemize}[leftmargin=2.4em,itemsep=0.3em]
    \item[\normsv] A foundational rule of the stateful protocol.
    Binding on every Phase~1 implementation. Revising a
    \normsv{} statement requires a separate alignment round and
    invalidates every \derived{} statement downstream of it.

    \item[\derived] Logically derived from one or more \normsv{}
    statements. Negotiable \emph{only} by showing the derivation
    is incorrect or that the rule it depends on has been revised.

    \item[\deferred] A gap deliberately left open. Each \deferred{}
    point names the artifact (companion paper, future RFC, or
    Appendix~B sub-section) that will resolve it.
\end{itemize}

Four typographic boxes appear:

\begin{itemize}[leftmargin=2.4em,itemsep=0.3em]
    \item A red \textbf{Invariant} box pins the binding commitment of
    a section in one sentence.
    \item A blue \textbf{Note} box adds clarification or pointers to
    related documents.
    \item An orange \textbf{Transitional misalignment} box names a
    specific place where the current \texttt{EasyNet-Cli}
    implementation contradicts the rule. These are debts, tracked
    explicitly, to be discharged in Phase~1.
    \item A purple \textbf{Open decision} box marks a choice that
    must be resolved before implementation proceeds.
\end{itemize}

\textbf{Reading order.} If reading for the first time, read
\S\ref{sec:state-object} through \S\ref{sec:receipt-index} in order;
the dependency graph is acyclic. \S\ref{sec:coverage} (the coverage
matrix) is best read \emph{after} the rest, with the companion
appendix open in a second window.

% ──────────────────────────────────────────────────────────────────────────
\section{State Object}
\label{sec:state-object}

\subsection{Definition}

\normsv{}
\begin{definition}[State object]
A \emph{state object} is a triple
\[
    \mathcal{S} = \bigl(\,t,\ k,\ v(\tau)\,\bigr)
\]
where $t$ is a \emph{state type}, $k$ is a \emph{state key} unique
within $t$, and $v(\tau)$ is the value of the object at logical time
$\tau$, evolving under a sequence of \emph{transitions}. The pair
$(t,k)$ is the object's identity; $v$ is its content.
\end{definition}

\subsection{State type}
\label{sec:state-type}

\normsv{}
A \emph{state type} $t$ is a stable string identifier with the
syntactic shape \texttt{<namespace>.<noun>}, both lowercase ASCII.
Examples: \texttt{easynet.page}, \texttt{easynet.web\_app}.
A state type names a \emph{class} of state objects sharing one
schema, one transition vocabulary, one canonical projection, and
one view declaration set. Two implementations claiming to host the
same state type \textbf{MUST} agree on all four byte-for-byte.

\subsection{State key}
\label{sec:state-key}

\normsv{}
A \emph{state key} $k$ is a tuple of typed scalars whose tuple
schema is fixed by the state type. The state key uniquely
identifies one instance within $t$. State key fields \textbf{MUST}
include enough information to derive the canonical authority for
the object (\S\ref{sec:authority}); placing the
\texttt{owner\_agent} URI inside the state key is the standard
mechanism, because authority that lived as a mutable field could
be rewritten by a transition, opening a privilege-escalation
primitive the protocol refuses to expose.

\subsection{Lifecycle}

\derived{} A state object's lifecycle is the sequence
\[
    \emptyset \;\xrightarrow{\;\text{first canonical transition}\;}\; v_1
    \;\to\; v_2 \;\to\; \cdots \;\to\; v_n \;\to\; \text{(terminal)}
\]
Each $v_i$ is the canonical value of the object after the
$i$-th canonical transition; transitions are defined in
\S\ref{sec:transition}. The terminal state is type-specific (e.g.
\texttt{Deleted} for many object types); after the terminal,
further transitions \textbf{MUST} be rejected.

\subsection{Identity immutability}
\label{sec:identity-immutability}

\begin{invariant}[SO-INV-1]
The state key of a state object \textbf{MUST NOT} change after the
first \texttt{CanonicalReceipt} for that key has been emitted.
A state object cannot be renamed, re-keyed, or moved across types;
the only way to ``rename'' is to delete one object and create a
distinct one.
\end{invariant}

\textbf{Rationale.} The audit chain in \S\ref{sec:receipt-index}
indexes receipts by state key. A mutable key would break receipt
enumeration: receipts emitted before the key change would no
longer be discoverable through the new key. Implementations are
expected to enforce SO-INV-1 at the receipt-validation layer (a
canonical transition that would mutate any state-key field is
rejected before apply).

\subsection{Relationship to abilities}

\derived{} A state object is \emph{not} an ability. Abilities are
the verbs that mutate a state object; the object itself is the
noun. One ability \textbf{MAY} affect zero or more state objects
(per its declared transition schema, \S\ref{sec:transition}); one
state object is acted on by one or more abilities. The protocol
does not permit an ability to invent its own state object types
at runtime --- every transition references a state type known
statically through schema registration.

\subsection{Schema-level canonicality marking}
\label{sec:canonicality-marking}

\normsv{}
The schema of every state type \textbf{MUST} annotate every field
with an explicit Boolean flag \texttt{canonical}. Fields with
\texttt{canonical:true} participate in
\texttt{canonical\_state\_hash} (\S\ref{sec:canonical-hash}); fields
with \texttt{canonical:false} do not. There is no implicit default;
omitting the flag is a schema error.

\begin{notebox}
The Phase~1 implementation in \texttt{EasyNet-Cli} \emph{does not
yet} carry this annotation on \texttt{AbilityDescriptor} or on
state-object schemas. The schema upgrade and the validator that
refuses unannotated schemas are the first two work items of
Phase~1a. See \S\ref{sec:coverage} for the appendix invariants
this section anchors.
\end{notebox}

% ──────────────────────────────────────────────────────────────────────────
\section{Canonical State Hash}
\label{sec:canonical-hash}

\subsection{Three-step computation}

\normsv{}
The \texttt{canonical\_state\_hash} of a state object is computed in
three deterministic, sequentially-composed steps:

\begin{enumerate}[leftmargin=2.5em]
    \item \textbf{Canonical projection.} A pure function
    \[
        \pi : v(\tau) \;\longmapsto\; v_{\text{can}}(\tau)
    \]
    that selects the subset of fields with
    \texttt{canonical:true} (\S\ref{sec:canonicality-marking}) and
    discards everything else. The projection \textbf{MUST} be
    schema-driven; an implementation that hand-rolls the
    projection for a specific state type produces a hash
    incompatible with peer implementations.

    \item \textbf{Canonical serialization.} A pure function
    \[
        \sigma : v_{\text{can}} \;\longmapsto\; b \in \{0,1\}^*
    \]
    producing a deterministic byte string. The serialization
    \textbf{MUST} be RFC~8785 (JSON Canonicalization Scheme,
    JCS): every implementation hashing the same canonical
    projection produces the same bytes regardless of language,
    map iteration order, or floating-point printer.

    \item \textbf{Hashing.} A pure function
    \[
        h : b \;\longmapsto\; \texttt{canonical\_state\_hash}
    \]
    using a single fixed cryptographic hash (Phase~1: SHA-256;
    upgradeable via a schema-level hash version field, but not
    swappable per call).
\end{enumerate}

\begin{decisionbox}
\textbf{Open: future hash agility.} SHA-256 is the Phase~1 choice.
A migration path to an algorithm with stronger collision margin
(SHA3-256 or BLAKE3) is required eventually. Mechanism: a per-state-type
schema field \texttt{hash\_version} with a registry of algorithms.
The mechanism is specified in this RFC; the migration policy is
deferred to RFC-009.
\end{decisionbox}

\subsection{CSH-INV-1: hash over projection, not over raw state}

\begin{invariant}[CSH-INV-1]
\texttt{canonical\_state\_hash} \textbf{MUST} be computed over
$\sigma(\pi(v))$, not over $\sigma(v)$ or any other transformation
of $v$. An implementation that hashes raw state, or that performs
its own ad-hoc projection, produces values incompatible with
conforming peers and is non-conformant.
\end{invariant}

\textbf{Why this invariant exists.} Without an explicit projection
step, distinct implementations diverge on which fields enter the
hash even when their schemas agree --- different languages serialize
maps in different orders, different SDKs canonicalize floats
differently, and over time a field accidentally tagged
\texttt{canonical:false} in one SDK and \texttt{canonical:true} in
another produces silently different hashes. CSH-INV-1 forces the
projection to be a first-class step, schema-driven, auditable.

\subsection{CSH-INV-2: content-as-blob boundary}
\label{sec:content-as-blob}

\begin{invariant}[CSH-INV-2]
A state object field whose semantically-meaningful value is large
(e.g. an HTML body, an image, a build artifact) \textbf{MUST} be
represented in canonical state as a content hash (a fixed-length
byte string), with the actual bytes stored in a content-addressed
blob store external to the state object. Embedding the bytes
directly in canonical state is non-conformant.
\end{invariant}

\begin{notebox}
The blob-store interface (\texttt{put}/\texttt{get}/\texttt{exists},
ref-counting and garbage-collection rules) is \deferred{} to
RFC-007. Phase~1 may implement a minimal local blob store that
satisfies \texttt{put}/\texttt{get}/\texttt{exists}; ref-counting
and GC may be stubbed.
\end{notebox}

\subsection{Hash and version binding}

\derived{} The pair $(\texttt{canonical\_state\_hash},\,
\texttt{state\_version})$ together identifies one specific value
in the lifecycle of one state object. \texttt{state\_version} is a
monotonic integer equal to the number of \texttt{CanonicalReceipt}s
that have advanced the object's state key
(\S\ref{sec:receipt-classes}). A verifier replaying canonical
receipts produces a sequence of pairs with strictly increasing
versions; if two distinct hashes are observed at the same version,
the receipt log is corrupt.

% ──────────────────────────────────────────────────────────────────────────
\section{Receipt Classes}
\label{sec:receipt-classes}

\subsection{CanonicalReceipt}

\normsv{}
A \texttt{CanonicalReceipt} witnesses one canonical transition of
one state object. Its required fields are:

\begin{lstlisting}
CanonicalReceipt {
    receipt_class:       "canonical",
    state_type:          string,
    state_key:           tuple,
    pre_state_hash:      bytes(32),     // canonical_state_hash
                                        // BEFORE the transition;
                                        // null/zero only for the
                                        // creation transition
    post_state_hash:     bytes(32),     // canonical_state_hash
                                        // AFTER the transition
    state_version:       uint64,
    transition_id:       string,        // §4
    attempt_id:          string,        // §3.4
    owner_signature:     bytes,         // §5.2
    observed_at:         timestamp
}
\end{lstlisting}

A \texttt{CanonicalReceipt} \textbf{MUST} be replayable offline:
given a sequence of canonical receipts in version order, a verifier
reconstructs the same canonical state without consulting any
network channel.

\subsection{OperationalReceipt}

\normsv{}
An \texttt{OperationalReceipt} records a runtime or observability
event. It \textbf{MUST NOT} carry \texttt{pre\_state\_hash} or
\texttt{post\_state\_hash} that refer to canonical state; the
canonical chain is closed under \texttt{CanonicalReceipt}s alone.
An \texttt{OperationalReceipt} \textbf{MAY} carry runtime metadata
(executor node, pid, port, log references, progress fractions),
\textbf{MAY} carry the \texttt{caller\_context} channel
(\S\ref{sec:caller-context}), and is not required for replay
correctness: the loss of any subset of operational receipts
\textbf{MUST NOT} affect the recoverability of canonical state.

\subsection{No mixing}

\normsv{}
A single ability invocation \textbf{MUST NOT} produce both a
\texttt{CanonicalReceipt} and an \texttt{OperationalReceipt} for
the same logical effect. A long-running canonical transition
produces many operational progress receipts \emph{plus} exactly
one terminal canonical receipt; this is one logical effect spread
across receipts of different classes, but each receipt witnesses a
distinct event (a progress observation versus the canonical
commit) and the classification is unambiguous.

\subsection{RC-INV-1: receipt binding}
\label{sec:receipt-binding}

\begin{invariant}[RC-INV-1]
Every receipt \textbf{MUST} carry the triple
\[
    \bigl(\texttt{state\_key},\ \texttt{transition\_id},\ \texttt{attempt\_id}\bigr)
\]
identifying exactly which transition attempt the receipt belongs
to. \texttt{attempt\_id} \textbf{MUST} be unique within a
\texttt{(state\_key,\ transition\_id)} pair and \textbf{MUST} be
preserved across all operational and canonical receipts
participating in the same attempt.
\end{invariant}

\textbf{Why this invariant exists.} Without \texttt{attempt\_id},
two concurrent invocations of the same transition on the same
state key cannot be told apart in the receipt log: their progress
streams interleave indistinguishably, and a verifier cannot
attribute a given progress receipt to a specific canonical
terminal. With \texttt{attempt\_id}, every receipt is unambiguously
attributable.

\subsection{Long-running transition pattern}
\label{sec:long-running}

\derived{} A long-running transition (one whose wall-clock duration
spans more than a single network round-trip) emits a receipt
sequence of the form

\begin{lstlisting}
on success:
    [1] OperationalReceipt   transition.started
    [2] OperationalReceipt   transition.progress     (0..N, lossy)
    [3] CanonicalReceipt     transition.completed
                              pre_state_hash  = hash before attempt
                              post_state_hash = hash after success

on failure:
    [1] OperationalReceipt   transition.started
    [2] OperationalReceipt   transition.progress     (0..N, lossy)
    [3] CanonicalReceipt     transition.failed
                              pre_state_hash  = hash before attempt
                              post_state_hash = hash after failure
\end{lstlisting}

The properties this implies are sharp:
\begin{itemize}[leftmargin=2em,itemsep=0.3em]
    \item Canonical state advances at exactly one point per attempt,
    namely the terminal canonical receipt.
    \item Replay from \texttt{CanonicalReceipt}s alone reaches the
    correct terminal; the transient ``in-flight'' state is invisible
    to a verifier without operational receipts.
    \item Progress receipts are observability, not history. A peer
    connecting late receives only the canonical terminal; one
    connecting early sees the progress stream as well.
\end{itemize}

\subsection{Caller-context channel}
\label{sec:caller-context}

\normsv{}
An \texttt{OperationalReceipt} \textbf{MAY} carry a structured
\texttt{caller\_context} map alongside (not replacing)
\texttt{envelope.caller}. The map records observations the call
site made about the request; its presence \textbf{MUST NOT}
influence \texttt{canonical\_state\_hash}, and its absence
\textbf{MUST NOT} affect canonical replay correctness. Specific
field schema is delegated to the call site (\S\ref{sec:caller-uri}
specifies one in the principal-URI context). RFC-006 main commits
only to the channel's existence and to its non-canonical nature.

\subsection{Concurrent transition admissibility}
\label{sec:concurrency}

\begin{invariant}[RC-INV-2]
Every canonical transition attempt \textbf{MUST} declare its
expected \texttt{pre\_state\_hash}. A canonical receipt whose
declared \texttt{pre\_state\_hash} does not match the current
\texttt{canonical\_state\_hash} of its state key \textbf{MUST}
be rejected at receipt-validation time, before any apply step.
\end{invariant}

This is optimistic concurrency: two attempts on the same state key
may proceed in parallel; whichever's terminal canonical receipt
arrives first wins, the other's receipt is rejected (the loser
re-reads state and retries). RFC-006 main \emph{commits} to
optimistic concurrency rather than strict serialization because the
former preserves liveness under network partition: a node that
cannot temporarily reach the owner of a state key can still queue
attempts whose viability is checked at the owner's next contact.
Strict serialization, while simpler to implement, requires a
single in-flight transition per key globally and is
network-partition-fragile.

\begin{notebox}
RC-INV-2's \texttt{pre\_state\_hash} field is the same field
already required in \S\ref{sec:receipt-classes}. The novelty here
is the validation rule: rejection is at the receipt log's gate,
not at the apply step. An apply-time check would let two attempts
both complete locally and only collide at convergence time, which
is not an admissible failure mode.
\end{notebox}

% ──────────────────────────────────────────────────────────────────────────
\section{Transition Ability}
\label{sec:transition}

\subsection{transition\_id}

\normsv{}
A \emph{transition} is an ability whose schema declares it as
canonical or operational. Every transition is identified by a
stable \texttt{transition\_id} string of the form

\[
    \texttt{<ability-qualified-name>@v<version>}
\]

where \texttt{<ability-qualified-name>} follows the ordinary EasyNet
ability naming scheme (\texttt{<namespace>.<verb>}) and
\texttt{<version>} is a monotonically-increasing integer assigned by
the schema registry when the ability's transition schema changes in
a non-backward-compatible way.

Examples: \texttt{webapp.build@v1}, \texttt{pages.publish@v2}.

\textbf{Schema evolution.} A receipt referring to
\texttt{webapp.build@v1} is forever validated under the v1
schema; a transition schema upgrade allocates a new
\texttt{transition\_id} (\texttt{@v2}). Old receipts replay
under their original schema; new attempts use the new
\texttt{transition\_id}. The receipt log is therefore stable
under unbounded schema evolution; this is the same discipline an
append-only event store applies to schema versioning.

\subsection{Transition schema fields}

\normsv{}
A transition schema declares, at minimum:

\begin{lstlisting}
{
    transition_id:    string,                 // see §4.1
    transition_class: "canonical" | "operational",
    pre_state:        set<state_value>,       // admissible inputs
    post_state:       set<state_value>,       // admissible outputs
    args_schema:      JSONSchema,
    receipt_schema:   JSONSchema              // schema of the
                                              // receipt this
                                              // transition emits
}
\end{lstlisting}

The \texttt{pre\_state} and \texttt{post\_state} sets are interpreted
relative to the state type's state space; an admissible attempt is
one whose owning object's current state value is in
\texttt{pre\_state}, and whose successful outcome lies in
\texttt{post\_state}.

\subsection{TA-INV-1: receipt-to-transition binding}

\begin{invariant}[TA-INV-1]
Every receipt \textbf{MUST} reference its transition's
\texttt{transition\_id}. Validators \textbf{MUST} resolve that
\texttt{transition\_id} against the schema registry and validate
the receipt's payload against the schema's
\texttt{receipt\_schema}. Receipts referencing an unknown
\texttt{transition\_id} \textbf{MUST} be rejected.
\end{invariant}

This is the schema-validation hook for the entire receipt log.
Without TA-INV-1, replay must trust that historic receipts
conformed to whatever schema was in force at the time, with no way
to verify; with TA-INV-1, the schema registry is the source of
truth, and the receipt log is auditable against it indefinitely.

\subsection{Validation timing}

\normsv{}
Transition schema validation \textbf{MUST} occur at
receipt-validation time (the moment a receipt is admitted to the
log), not at apply time. A receipt that fails schema validation
\textbf{MUST} be rejected before any state mutation; an
implementation that allows non-conformant receipts to apply
locally and only rejects them at convergence time is
non-conformant.

\subsection{Cross-state-space composition}
\label{sec:cross-state-space}

\normsv{}
A state type \textbf{MAY} declare multiple parallel state spaces
on the same state key (the canonical / runtime overlay split of
Appendix~B is the minimal interesting case; types declaring
review, billing, or audit overlays are admissible future cases).
When a state type declares $n \geq 2$ state spaces, its schema
\textbf{MUST} include a \emph{cross-product admissibility matrix}
specifying which combinations of values across the spaces are
allowed; transitions producing an inadmissible combination
\textbf{MUST} be rejected at receipt-validation time.

\begin{notebox}
The cross-product matrix is what prevents a runtime crash from
being modelled as a canonical state change. With only one state
space, an implementation cannot express the difference; with two
spaces and an explicit admissibility matrix, the impossibility of
\texttt{(canonical=Building, runtime=Running)} is a schema-level
fact, not a runtime invariant the application must remember.
\end{notebox}

% ──────────────────────────────────────────────────────────────────────────
\section{Authority Domain and Delegation}
\label{sec:authority}

\subsection{Owner agent}

\normsv{}
Every state object has exactly one \emph{owner agent}, identified
by an EasyNet agent URI. The owner agent is the \emph{canonical
authority} for the object: every \texttt{CanonicalReceipt} for the
object \textbf{MUST} carry the owner agent's signature
(\S\ref{sec:owner-sig}).

The owner agent's identity \textbf{MUST} be embedded in the
state key (\S\ref{sec:state-key}), not in a mutable field.
Embedding it in the key turns ``who can sign canonical
transitions for this object'' into a property of the object's
identity rather than a property of its content. Because keys are
immutable (SO-INV-1), authority cannot be silently rewritten by a
transition; the only way to change ownership is to delete the
object and create a distinct one with a different key.

\subsection{Owner signature on canonical receipts}
\label{sec:owner-sig}

\normsv{}
A \texttt{CanonicalReceipt} for state key $k$ \textbf{MUST} carry
a cryptographic signature
\[
    \texttt{owner\_signature} = \mathrm{Sign}_{\mathrm{sk}_{\mathrm{owner}(k)}}\!\bigl(\sigma(\text{receipt body})\bigr)
\]
where $\mathrm{sk}_{\mathrm{owner}(k)}$ is the private key bound to
the owner agent identity in $k$, and $\sigma$ is the canonical
serialization of \S\ref{sec:canonical-hash}. A canonical receipt
without owner signature \textbf{MUST} be rejected; a receipt whose
signature does not verify against the owner agent's published key
\textbf{MUST} be rejected; verification \textbf{MUST} be performed
at receipt-validation time.

\begin{decisionbox}
\textbf{Open: signature scheme.} Phase~1 commits to Ed25519 for
its small key size and well-understood security profile.
Threshold or multi-signature schemes (for shared-ownership
objects, e.g. organisations) are \deferred{} to RFC-008.
\end{decisionbox}

\subsection{Delegated execution}

\normsv{}
A canonical transition \textbf{MAY} be \emph{executed} by a node
other than the owner --- the source of a webapp may live on
agent~$A$, the build farm on agent~$B$, the serve runtime on
agent~$C$ --- but the owner remains the canonical authority. The
delegated execution flow is:

\begin{enumerate}[leftmargin=2.5em,itemsep=0.3em]
    \item The owner emits an authorisation token (or signs an
    attempt-prefix) addressed to the delegate agent and bound to
    \texttt{(state\_key, transition\_id, attempt\_id)}.
    \item The delegate executes the transition; on success or
    failure, it constructs the canonical receipt body but does
    \emph{not} sign it.
    \item The delegate returns the receipt body to the owner.
    \item The owner verifies the body, signs it with
    $\mathrm{sk}_{\mathrm{owner}(k)}$, and emits the signed
    \texttt{CanonicalReceipt}.
\end{enumerate}

A delegate-emitted receipt without owner signature \textbf{MUST
NOT} advance canonical state. The owner's signature is the
authoritative act; the delegate's execution is mechanical labour.

\begin{notebox}
The authorisation token format and the round-trip latency
budget for the delegate $\to$ owner signing handshake are
\deferred{} to RFC-008. Phase~1 may implement a synchronous
loop in which the owner signs canonical receipts on demand from
co-located delegates.
\end{notebox}

\subsection{The principal/ caller URI scheme}
\label{sec:caller-uri}

\normsv{}
The \texttt{envelope.caller} field of every invocation \textbf{MUST}
be a URI under one of two namespaces:

\begin{itemize}[leftmargin=2em,itemsep=0.3em]
    \item \textbf{Agent caller.}
    $\texttt{easynet:///agent/<id>}$ --- the caller is an EasyNet
    agent. This is the original (pre-stateful) form.

    \item \textbf{Principal caller.}
    $\texttt{easynet:///principal/<class>/<\,\ldots\,>}$ --- the
    caller is a principal that is not an EasyNet agent. The
    \texttt{<class>} segment is one of a small set of reserved
    classes (\texttt{human-anon}, \texttt{human-auth},
    \texttt{bot-auth}; new classes require an RFC). Subsequent
    segments are class-specific.
\end{itemize}

The reserved class \texttt{human-anon} has the form

\[
    \texttt{easynet:///principal/human-anon/<ip-hash>/<session-id|"-">}
\]

where \texttt{ip-hash} is a stable hash (e.g.
\texttt{HMAC(local\_salt, ip).hex[..16]}) and \texttt{session-id}
is an opaque token or the literal \texttt{"-"} when absent. The
class \texttt{human-anon} is the conformant placeholder for a
hub-translated public browser request; using the hub's own agent
URI as caller for such requests is non-conformant.

\textbf{Disambiguation discipline.} Authorisation, scope-check, and
audit code paths \textbf{MUST} discriminate \texttt{agent/} from
\texttt{principal/} by URI prefix and \textbf{MUST NOT} treat a
\texttt{principal/} caller as carrying agent-level identity
guarantees. Specifically: a \texttt{principal/human-anon/...} caller
has no signature, no key, no realm membership; any code path that
relies on those for an \texttt{agent/} caller \textbf{MUST} fail
fast on a \texttt{principal/} caller rather than degrading
silently.

\begin{warningbox}
Phase~1 implementation: the existing \texttt{envelope.caller} parser
in \texttt{EasyNet-Cli} assumes \texttt{agent/} form throughout.
Acceptance of \texttt{principal/} URIs requires extending the parser,
extending scope-check, and adding a fail-fast classifier in
\texttt{ability\_proxy}. This is Phase~1c work.
\end{warningbox}

% ──────────────────────────────────────────────────────────────────────────
\section{View Projection}
\label{sec:view-projection}

\subsection{Definition}

\normsv{}
A \emph{view} is a pure projection of a state object's
$(\text{canonical state},\ \text{runtime overlay state})$ pair, and,
optionally, of the \texttt{caller\_context} of the request that
asked for the view. A view computation \textbf{MUST NOT} mutate any
state, \textbf{MUST NOT} consult side channels (network, time of
day, load, locale) for its output content, and \textbf{MUST}
produce identical bytes when given identical inputs.

(View \emph{caching} may consult time. View \emph{content} may not.)

\subsection{Multi-view as first-class}

\normsv{}
A state type \textbf{MAY} declare $n \geq 0$ named views in its
schema. Any state type intended for exposure through a Hub-class
deployment \textbf{MUST} declare at least the two reserved views

\begin{itemize}[leftmargin=2em,itemsep=0.2em]
    \item \texttt{human\_view} --- HTTP-renderable (HTML, image,
    stream).
    \item \texttt{agent\_view} --- structured, machine-actionable.
\end{itemize}

The two views \textbf{MUST} produce structurally distinct outputs;
serving the same byte sequence to a human GET and an agent GET on
the URA channel is non-conformant content negotiation, not
multi-view.

\subsection{agent\_view minimum content}
\label{sec:agent-view-min}

\normsv{}
The \texttt{agent\_view} of any state object \textbf{MUST}
include, at minimum:

\begin{enumerate}[leftmargin=2em,itemsep=0.25em]
    \item A snapshot of the canonical state (the values of all
    canonical fields at the current version).
    \item A snapshot of the runtime overlay state, when the
    state type declares one.
    \item An \texttt{ability\_surface} listing the
    \texttt{transition\_id}s admissible from the current state.
    \item A pointer to the receipt-history view of this state
    object (\S\ref{sec:receipt-index}).
\end{enumerate}

It \textbf{MAY} additionally include an \texttt{eal\_hint} (a
ready-to-paste EAL invocation that advances the state object) and
implementation-specific fields enabled by the state type's view
schema.

\subsection{VP-INV-1: ability\_surface is derived, not registry-snapshot}

\begin{invariant}[VP-INV-1]
The \texttt{ability\_surface} of an \texttt{agent\_view}
\textbf{MUST} be derived from the current state of the object and
the transition schema (specifically: the set of
\texttt{transition\_id}s whose \texttt{pre\_state} contains the
current state value, intersected with the cross-product
admissibility matrix of \S\ref{sec:cross-state-space}). It
\textbf{MUST NOT} be a static snapshot of the agent's ability
registry.
\end{invariant}

\textbf{Why this invariant exists.} An \texttt{ability\_surface}
that simply lists ``every ability the agent has ever registered''
defeats state-gating: a peer agent reading the view would attempt
calls that are inadmissible from the current state, and the
state-gating logic would have to be replicated at every caller.
Deriving the surface from current state and schema makes the view
a genuine affordance: the listed transitions are precisely those
that will succeed if attempted now.

\subsection{Determinism}

\normsv{}
Every view of a state type is a deterministic function of its
declared input set. The default input set is
$(\text{canonical},\ \text{runtime})$; views that consult
\texttt{caller\_context} \textbf{MUST} declare it as an explicit
input. No view \textbf{MAY} consult \texttt{caller\_context}
ambiently or partially; the dependency is either declared and
deterministic, or absent.

% ──────────────────────────────────────────────────────────────────────────
\section{Receipt Indexing by State Key}
\label{sec:receipt-index}

\subsection{Dual indexing requirement}

\normsv{}
The receipt log \textbf{MUST} support enumeration along two axes
simultaneously:

\begin{itemize}[leftmargin=2em,itemsep=0.25em]
    \item \textbf{By \texttt{receipt\_id}} --- preserves the
    pre-stateful ``one invocation is traceable end-to-end''
    property of the original EasyNet protocol.
    \item \textbf{By \texttt{state\_key}} --- enables the new
    ``one state object's entire life is traceable end-to-end''
    property.
\end{itemize}

A log that supports only one axis is non-conformant.

\subsection{URA mapping}

\normsv{}
Every state object \textbf{MUST} be addressable by a stable
URA derived from its state key alone, via a mapping function
declared by the state type's schema. The mapping
\textbf{MUST} be bijective on $(\texttt{state\_type},
\texttt{state\_key}) \leftrightarrow \texttt{ura}$ within a realm:
two state objects map to the same URA if and only if they have
the same key under the same type, and the URA can be inverted to
recover both.

The standard form is

\[
    \texttt{easynet:///r/<realm>/agent/<owner>/<state\_type>/<slug>}
\]

where \texttt{<slug>} is the rendering of the non-owner part of
the state key; the schema specifies the rendering function.

\subsection{Receipt-history enumeration via URA}

\normsv{}
Given a state object's URA, a verifier \textbf{MUST} be able to
enumerate the full ordered sequence of \texttt{CanonicalReceipt}s
that have advanced the object's state, up to a chosen
\texttt{state\_version} (or up to the present), without
consulting any channel other than the receipt log addressable
through that URA.

The receipt-history view \textbf{MAY} be exposed as a sub-URA of
the form $\texttt{<state\_object\_ura>/receipts}$, but the URA
form is not protocol-fixed; the binding requirement is that the
enumeration is reachable through the same trust boundary as the
state object itself.

\subsection{Append-only}

\normsv{}
The receipt log is append-only. Receipts \textbf{MUST NOT} be
deleted, renumbered, or rewritten; corrections are made by emitting
a new receipt that supersedes (but does not erase) the prior one.
Implementations that compact, archive, or shard the log
\textbf{MUST} preserve the externally-observable enumeration
contract: a client repeating an enumeration query receives the
same sequence (or a strict prefix), modulo new receipts appended
since.

% ──────────────────────────────────────────────────────────────────────────
\section{Invariant Coverage Matrix}
\label{sec:coverage}

This section verifies that the main body anchors every invariant
declared in Appendix~B (the companion pressure-test) and that
every locally-declared invariant in this document is anchored to
one main-body section. If any row in either table is missing or
inconsistent, the main body is incomplete and must be revised
before Phase~1 implementation begins.

\subsection{Appendix~B invariants $\rightarrow$ main-body anchors}

\begin{center}
\renewcommand{\arraystretch}{1.25}
\begin{tabularx}{\linewidth}{l X l}
\toprule
\textbf{Invariant} & \textbf{Statement (abridged)} & \textbf{Anchor} \\
\midrule
TR-INV-1  & Long-running transition completion semantics       & \S\ref{sec:long-running} \\
TR-INV-2  & Progress receipts not in canonical chain           & \S\ref{sec:receipt-classes}, \S\ref{sec:long-running} \\
TR-INV-3  & Schema-level canonicality marking                  & \S\ref{sec:canonicality-marking}, \S\ref{sec:canonical-hash} \\
TR-INV-4  & Owner authority over delegated execution           & \S\ref{sec:owner-sig} \\
TR-INV-5  & Runtime overlay observable but not canonical       & \S\ref{sec:receipt-classes}, \S\ref{sec:cross-state-space} \\
TR-INV-6  & View $\neq$ State                                  & \S\ref{sec:view-projection} \\
TR-INV-7  & Multi-view as first-class                          & \S\ref{sec:view-projection}, \S\ref{sec:agent-view-min} \\
TR-INV-8  & Cross-state-space composition                      & \S\ref{sec:cross-state-space} \\
TR-INV-9  & Concurrent canonical transitions                   & \S\ref{sec:concurrency} (RC-INV-2) \\
TR-INV-10 & Content-as-blob boundary                           & \S\ref{sec:content-as-blob} (CSH-INV-2) \\
TR-INV-11 & URA-receipt traceability                           & \S\ref{sec:receipt-index} \\
TR-INV-12 & Hub-translated caller identity                     & \S\ref{sec:caller-uri} \\
TR-INV-13 & Caller-context as observability channel            & \S\ref{sec:caller-context} \\
\bottomrule
\end{tabularx}
\end{center}

\subsection{Main-body invariants (self-contained)}

\begin{center}
\renewcommand{\arraystretch}{1.25}
\begin{tabularx}{\linewidth}{l X l}
\toprule
\textbf{Invariant} & \textbf{Statement (abridged)} & \textbf{Defined in} \\
\midrule
SO-INV-1  & State-key immutability after first canonical receipt & \S\ref{sec:identity-immutability} \\
CSH-INV-1 & Hash computed over projection, not raw state         & \S\ref{sec:canonical-hash} \\
CSH-INV-2 & Content-as-blob boundary                             & \S\ref{sec:content-as-blob} \\
RC-INV-1  & Receipt binding to (state\_key, transition\_id, attempt\_id) & \S\ref{sec:receipt-binding} \\
RC-INV-2  & Pre-state-hash declaration and rejection             & \S\ref{sec:concurrency} \\
TA-INV-1  & Receipt-to-transition binding                        & \S\ref{sec:transition} \\
VP-INV-1  & ability\_surface derived from state and schema       & \S\ref{sec:view-projection} \\
\bottomrule
\end{tabularx}
\end{center}

\subsection{Coverage statement}

\derived{} The two tables above account for every TR-INV-*
invariant declared in the companion appendix and for every
SO/CSH/RC/TA/VP-INV-* invariant declared in this main body. No
invariant is left without a normative anchor; no anchor is
referenced by zero invariants. The Phase~1 implementation work
plan
(\S\ref{sec:phase1}) is therefore well-founded: each work item
discharges a specific subset of the rows above.

% ──────────────────────────────────────────────────────────────────────────
\section{Phase 1 implementation work plan}
\label{sec:phase1}

\derived{} The work this RFC immediately authorises is bounded.
Each item below is a Phase~1 deliverable in
\texttt{EasyNet-Cli}, mapping to one or more invariants from
\S\ref{sec:coverage}.

\begin{enumerate}[leftmargin=2.5em,itemsep=0.3em]
    \item \textbf{1a. Schema-level \texttt{canonical:bool}
    annotation.} Extend \texttt{AbilityDescriptor} and state-type
    schemas with the \texttt{canonical} flag; reject schemas
    missing it. Discharges TR-INV-3, SO-INV-1 (key-field
    annotation).

    \item \textbf{1b. Two-class receipt split.} Introduce
    \texttt{CanonicalReceipt} and \texttt{OperationalReceipt}
    types, with the fields of \S\ref{sec:receipt-classes};
    \texttt{owner\_signature} stub-implemented (signing scheme
    choice from \S\ref{sec:owner-sig}). Discharges TR-INV-1,
    TR-INV-2, TR-INV-5, RC-INV-1, RC-INV-2, TA-INV-1.

    \item \textbf{1c. \texttt{principal/} caller URI parser and
    classifier.} Extend \texttt{envelope.caller} parser; add a
    fail-fast classifier; ensure scope-check distinguishes
    \texttt{agent/} from \texttt{principal/}. Discharges TR-INV-12.

    \item \textbf{1d. Minimal blob store.} A local content-addressed
    blob store implementing \texttt{put}/\texttt{get}/\texttt{exists};
    GC and ref-counting stubbed (resolved in RFC-007). Discharges
    TR-INV-10, CSH-INV-2.

    \item \textbf{1e. State-key receipt index.} Index the receipt
    log by state key in addition to receipt id; expose URA-keyed
    enumeration. Discharges TR-INV-11.
\end{enumerate}

The first reference state machine atop these primitives is
\texttt{easynet.page} (Appendix~A of the companion document, to be
drafted). The pressure-test machine is \texttt{easynet.web\_app}
(Appendix~B of the companion document, already drafted). View
projection (\S\ref{sec:view-projection}) and
\texttt{caller\_context} (\S\ref{sec:caller-context}) are exercised
end-to-end only when the Hub-class deployment is built; the
protocol surface is complete before that engineering work begins.

% ──────────────────────────────────────────────────────────────────────────
\section{Closing}
\label{sec:closing}

The intent of this RFC is not to specify every detail of a
stateful agent network; it is to specify the smallest surface a
conformant implementation must implement before further work can
proceed. Where the appendix pressure-tests revealed an invariant,
the main body addresses it; where the implementation work plan
requires a primitive, the main body defines it. Everything else
--- formal theorems, hub deployment, web reference machines,
cross-binding multi-realm consistency --- is deferred to companion
artefacts.

\paragraph{Companion artefacts.}
\begin{itemize}[leftmargin=2em,itemsep=0.2em]
    \item \texttt{AXON-RFC-006-stateful-easynet.md} ---
    Appendix~B (\texttt{easynet.web\_app}) and the thirteen
    TR-INV-* invariants this main body satisfies.
    \item Appendix~A (\texttt{easynet.page}, the simple worked
    example) --- \deferred{} to a follow-up commit on the
    companion document.
    \item RFC-007 --- blob store interface and GC rules.
    \item RFC-008 --- delegated execution authorisation tokens
    and threshold-signature ownership.
\end{itemize}

\paragraph{Bindingness.}
Sections marked \normsv{} are binding on every Phase~1 (and
later) implementation. Sections marked \derived{} are negotiable
only by showing the derivation incorrect. Sections marked
\deferred{} name the artefact that resolves them; an
implementation that needs the resolution sooner than that
artefact lands must request an out-of-band exception.

\vfill
\hrule
\vspace{0.4em}
\noindent\small
\textit{This document is a pre-implementation minimal normative
core. It is paired with the companion appendix (Appendix~B) which
served as its pressure test. If a future revision of either
document creates a conflict between them, the appendix wins on
matters of invariant coverage and the main body wins on matters
of normative form; reconciliation is a separate alignment round.}

\end{document}
